\documentclass[11pt,a4paper]{article}
\usepackage[T1]{fontenc}
\usepackage[utf8]{inputenc}
\usepackage{amsmath,amssymb,amsthm}
\usepackage[margin=1in]{geometry}
\usepackage[colorlinks=true,linkcolor=blue,urlcolor=blue]{hyperref}
\theoremstyle{plain}
\newtheorem{theorem}{Theorem}
\newtheorem{lemma}[theorem]{Lemma}
\newtheorem{proposition}[theorem]{Proposition}
\newtheorem{corollary}[theorem]{Corollary}
\theoremstyle{definition}
\newtheorem{remark}[theorem]{Remark}

\title{The empirical closed form for OEIS A227255, proved}
\author{Adrian Perez Fontelles\\ \small Independent researcher}
\date{1 September 2026}
\begin{document}
\maketitle

\begin{abstract}
OEIS A227255 carries an empirical closed form for $a(n)$ --- not a recurrence relating
consecutive terms but an explicit formula. It is true, and it is decidable rather than
empirical. The entry counts arrays subject to a condition on a bounded window of consecutive lines, so the lines of an array are the
vertices of a finite digraph, an array is a walk in it, and $a(n)$ is a walk count on
$S=1247$ vertices. Any function of the form $\sum_j p_j(n)\lambda_j^{\,n}$ satisfies the
monic linear recurrence whose characteristic polynomial is
$\prod_j(x-\lambda_j)^{\deg p_j+1}$; here that polynomial is $q(x)=(x-1)^{48}$, of degree
$48$. Two things are then checked exactly: that the walk count satisfies the same
recurrence from some index on, which the Cayley--Hamilton theorem reduces to a finite
computation, and that the two sequences agree at $48$ consecutive indices past that
point. A monic recurrence determines every later term from its predecessors, so agreement
there is agreement for good.
\end{abstract}

\noindent\small 2020 Mathematics Subject Classification. 05A15, 11B37, 15A18.\normalsize

\section{The conjecture}

OEIS A227255 is ``Number of n X 5 binary arrays indicating whether each 2 X 2 subblock of a larger binary array has lexicographically nondecreasing rows and columns, for some larger (n+1) X 6 binary array with rows and columns of the latter in lexicographically nondecreasing order.''. It has offset $1$ and begins
\[
11, 53, 450, 4578, 44379, 385212, 2925969, 19641271,\ \dots
\]
The entry states, as a conjecture contributed by its author and never marked settled:
\begin{quote}
Empirical polynomial of degree 47 (see link above).

Empirical: a(n) = (1/1010247037153000705636579512318796796754678251520000000000)*n\^{}47 - (1/2686827226470746557544094447656374459453931520000000000)*n\^{}46 + (1/7026681206853864812145392475073878939402240000000000)*n\^{}45 - (127/3664896472594368705942498820332650584080384000000000)*n\^{}44 + (72329/9439884853652161818336739385705312110510080000000000)*n\^{}43 - (81539/59872419790605677494736190607010436218880000000000)*n\^{}42 + (429055931/1975789853089987357326294290031344395223040000000000)*n\^{}41 - (16964821/566941134315634822762207830712007000064000000000)*n\^{}40 + (1827234659/491734657314581223824363934801230561280000000000)*n\^{}39 - (507972146243/1235640933764845126533017066936425512960000000000)*n\^{}38 + (1640156111807/39742837050916071321237391041814855680000000000)*n\^{}37 - (90602840371/24167941449881394722374089147049574400000000)*n\^{}36 + (6896619141543191/22177640389302926921707987687880785920000000000)*n\^{}35 - (386228642185150949/16392168983397815550827643073651015680000000000)*n\^{}34 + (472091170064351179/288021303757180869113090749193256960000000000)*n\^{}33 - (426965885290472681/4073028537980335522811384332025856000000000)*n\^{}32 + (47100747825947143625251/7616563366023227427657288700888350720000000000)*n\^{}31 - (41401107169380409667033/122847796226181087542859495175618560000000000)*n\^{}30 + (714437425832081190815995163/42013946309353931939657947350061547520000000000)*n\^{}29 - (589442711942561863039637/741052055901824357344703189876736000000000)*n\^{}28 + (4138300597776596584700524127/119821988437358893268404827806760960000000000)*n\^{}27 - (569185815975965973657984885269/408623704158172636017893387135877120000000000)*n\^{}26 + (163544534518752339012529560528817/3132781731879323542803849301375057920000000000)*n\^{}25 - (200994974625869228586313218455477/110568767007505536804841740048531456000000000)*n\^{}24 + (21290867886418332771208338997711489/362189192322016654652223684052254720000000000)*n\^{}23 - (232276730736190008644249112631866701/131705160844369692600808612382638080000000000)*n\^{}22 + (79265259917337030447380680303974383/1616916037151860065411712875233280000000000)*n\^{}21 - (60857585913006825102371224281427606751/48291892309602220620296491206967296000000000)*n\^{}20 + (2250002445476983293023101885534802801/75274384942512921288516847534080000000000)*n\^{}19 - (15024586569408678096426533515266502089541/23033961792408953914286155345428480000000000)*n\^{}18 + (1064011077877873794147065038577437650966911/81535442607731348976867601793679360000000000)*n\^{}17 - (26425325228994326799914489489167889830135313/110888201946514634608539938439403929600000000)*n\^{}16 + (831175224890362291169424698077063846965271259/210390561728878212984953008199761920000000000)*n\^{}15 - (247729484242257766360628691978964700287651283/4195823168382186583745501588029440000000000)*n\^{}14 + (112736477575864215276229064569328805498044789/143039426194847269900414826864640000000000)*n\^{}13 - (27426320212410665458099601016474314382917/2956199874857986926046445568000000000)*n\^{}12 + (13720116479138178578504440093160196340240816328873/145248590666078799957765678079549440000000000)*n\^{}11 - (1443427524169251064911674070485645151658026586043/1793192477358997530342786149130240000000000)*n\^{}10 + (325130524317293549697070523931430545850702704281/60671925925680367567989005045760000000000)*n\^{}9 - (157215426146907042983743454699228680841534721719/7032127362192147177814847643648000000000)*n\^{}8 - (57730612875783149128815834452990004774547179474761/2021824518222269715461491383144345600000000)*n\^{}7 + (1531260273681289171755604278348670505406567629149/976726820397231746599754291374080000000)*n\^{}6 - (11492431235654045096559324539290390817365777925329/689987732409187283847969281549260800000)*n\^{}5 + (1237685571728360417241018342162175103490344421103/10952186228717258473777290183321600000)*n\^{}4 - (12040107032687443295289959544247851472622437/22272823053760375542617851680000)*n\^{}3 + (4107286385195229860769185766558085320773/2292483225343906906575264000)*n\^{}2 - (413702984217720271054500548227873/110680160865928453800)*n + 3725941823309 for n>18
\end{quote}
As of the ``Last modified'' line on the live entry (Sep 10 2023 21:01:11, revision 8) this is still
recorded as empirical, and nothing on the entry records it as proved. Write
\[
f(n)\;=\;\frac{n^{47}}{1010247037153000705636579512318796796754678251520000000000} - \frac{n^{46}}{2686827226470746557544094447656374459453931520000000000} + \frac{n^{45}}{7026681206853864812145392475073878939402240000000000} - \frac{127 n^{44}}{3664896472594368705942498820332650584080384000000000} + \frac{72329 n^{43}}{9439884853652161818336739385705312110510080000000000} - \frac{81539 n^{42}}{59872419790605677494736190607010436218880000000000} + \frac{429055931 n^{41}}{1975789853089987357326294290031344395223040000000000} - \frac{16964821 n^{40}}{566941134315634822762207830712007000064000000000} + \frac{1827234659 n^{39}}{491734657314581223824363934801230561280000000000} - \frac{507972146243 n^{38}}{1235640933764845126533017066936425512960000000000} + \frac{1640156111807 n^{37}}{39742837050916071321237391041814855680000000000} - \frac{90602840371 n^{36}}{24167941449881394722374089147049574400000000} + \frac{6896619141543191 n^{35}}{22177640389302926921707987687880785920000000000} - \frac{386228642185150949 n^{34}}{16392168983397815550827643073651015680000000000} + \frac{472091170064351179 n^{33}}{288021303757180869113090749193256960000000000} - \frac{426965885290472681 n^{32}}{4073028537980335522811384332025856000000000} + \frac{47100747825947143625251 n^{31}}{7616563366023227427657288700888350720000000000} - \frac{41401107169380409667033 n^{30}}{122847796226181087542859495175618560000000000} + \frac{714437425832081190815995163 n^{29}}{42013946309353931939657947350061547520000000000} - \frac{589442711942561863039637 n^{28}}{741052055901824357344703189876736000000000} + \frac{4138300597776596584700524127 n^{27}}{119821988437358893268404827806760960000000000} - \frac{569185815975965973657984885269 n^{26}}{408623704158172636017893387135877120000000000} + \frac{163544534518752339012529560528817 n^{25}}{3132781731879323542803849301375057920000000000} - \frac{200994974625869228586313218455477 n^{24}}{110568767007505536804841740048531456000000000} + \frac{21290867886418332771208338997711489 n^{23}}{362189192322016654652223684052254720000000000} - \frac{232276730736190008644249112631866701 n^{22}}{131705160844369692600808612382638080000000000} + \frac{79265259917337030447380680303974383 n^{21}}{1616916037151860065411712875233280000000000} - \frac{60857585913006825102371224281427606751 n^{20}}{48291892309602220620296491206967296000000000} + \frac{2250002445476983293023101885534802801 n^{19}}{75274384942512921288516847534080000000000} - \frac{15024586569408678096426533515266502089541 n^{18}}{23033961792408953914286155345428480000000000} + \frac{1064011077877873794147065038577437650966911 n^{17}}{81535442607731348976867601793679360000000000} - \frac{26425325228994326799914489489167889830135313 n^{16}}{110888201946514634608539938439403929600000000} + \frac{831175224890362291169424698077063846965271259 n^{15}}{210390561728878212984953008199761920000000000} - \frac{247729484242257766360628691978964700287651283 n^{14}}{4195823168382186583745501588029440000000000} + \frac{112736477575864215276229064569328805498044789 n^{13}}{143039426194847269900414826864640000000000} - \frac{27426320212410665458099601016474314382917 n^{12}}{2956199874857986926046445568000000000} + \frac{13720116479138178578504440093160196340240816328873 n^{11}}{145248590666078799957765678079549440000000000} - \frac{1443427524169251064911674070485645151658026586043 n^{10}}{1793192477358997530342786149130240000000000} + \frac{325130524317293549697070523931430545850702704281 n^{9}}{60671925925680367567989005045760000000000} - \frac{157215426146907042983743454699228680841534721719 n^{8}}{7032127362192147177814847643648000000000} - \frac{57730612875783149128815834452990004774547179474761 n^{7}}{2021824518222269715461491383144345600000000} + \frac{1531260273681289171755604278348670505406567629149 n^{6}}{976726820397231746599754291374080000000} - \frac{11492431235654045096559324539290390817365777925329 n^{5}}{689987732409187283847969281549260800000} + \frac{1237685571728360417241018342162175103490344421103 n^{4}}{10952186228717258473777290183321600000} - \frac{12040107032687443295289959544247851472622437 n^{3}}{22272823053760375542617851680000} + \frac{4107286385195229860769185766558085320773 n^{2}}{2292483225343906906575264000} - \frac{413702984217720271054500548227873 n}{110680160865928453800} + 3725941823309.
\]

\section{The count is a walk count}

The entry's defining condition is local: it is a condition on a bounded window of consecutive lines. Take as vertices the
admissible configurations of such a window and put an edge from $u$ to $v$ when $v$ may
follow $u$, which the condition decides by inspection. An admissible array is then exactly a
walk, so with $M$ the adjacency matrix of that digraph on $S=1247$ vertices, and $u,w$ the
vectors recording which windows may start and end an array,
\[
a(n)\;=\;u^{\!\top}M^{\,g(n)}w
\]
for the linear function $g$ that the entry's shape supplies. In particular $a$ is
$C$-finite: it satisfies some monic integer linear recurrence, though not necessarily a short
one.

\section{A closed form is a recurrence}

\begin{lemma}\label{lem:ann}
Let $f(m)=\sum_{j}p_j(m)\lambda_j^{\,m}$ with the $\lambda_j$ distinct and the $p_j$
polynomials, and put $q(x)=\prod_j(x-\lambda_j)^{\deg p_j+1}$. Then $q$ applied to the shift
operator annihilates $f$: writing $q(x)=x^{r}-\sum_{i=1}^{r}c_ix^{r-i}$,
\[
f(m)\;=\;\sum_{i=1}^{r}c_i\,f(m-i)\qquad\text{for every }m.
\]
\end{lemma}

\begin{proof}
The shift operator $E$ acts on $m\mapsto\lambda^{m}p(m)$ as
$(E-\lambda)\bigl(\lambda^{m}p(m)\bigr)=\lambda^{m+1}\bigl(p(m+1)-p(m)\bigr)$, and
$p(m+1)-p(m)$ has degree one less than $p$. So $(E-\lambda)^{\deg p+1}$ kills
$\lambda^{m}p(m)$, and the factors of $q$ commute.
\end{proof}

Here $q(x)=(x-1)^{48}$, of degree $48$; the closed form is a polynomial in $n$, so this is $(x-1)^{48}$, and the recurrence it encodes is
\begin{equation}\label{eq:rec}
a(n)\;=\;48\,a(n-1) - 1128\,a(n-2) + 17296\,a(n-3) - 194580\,a(n-4) + 1712304\,a(n-5) - 12271512\,a(n-6) + 73629072\,a(n-7) - 377348994\,a(n-8) + 1677106640\,a(n-9) - 6540715896\,a(n-10) + 22595200368\,a(n-11) - 69668534468\,a(n-12) + 192928249296\,a(n-13) - 482320623240\,a(n-14) + 1093260079344\,a(n-15) - 2254848913647\,a(n-16) + 4244421484512\,a(n-17) - 7309837001104\,a(n-18) + 11541847896480\,a(n-19) - 16735679449896\,a(n-20) + 22314239266528\,a(n-21) - 27385657281648\,a(n-22) + 30957699535776\,a(n-23) - 32247603683100\,a(n-24) + 30957699535776\,a(n-25) - 27385657281648\,a(n-26) + 22314239266528\,a(n-27) - 16735679449896\,a(n-28) + 11541847896480\,a(n-29) - 7309837001104\,a(n-30) + 4244421484512\,a(n-31) - 2254848913647\,a(n-32) + 1093260079344\,a(n-33) - 482320623240\,a(n-34) + 192928249296\,a(n-35) - 69668534468\,a(n-36) + 22595200368\,a(n-37) - 6540715896\,a(n-38) + 1677106640\,a(n-39) - 377348994\,a(n-40) + 73629072\,a(n-41) - 12271512\,a(n-42) + 1712304\,a(n-43) - 194580\,a(n-44) + 17296\,a(n-45) - 1128\,a(n-46) + 48\,a(n-47) - a(n-48).
\end{equation}

\section{The proof}

\begin{lemma}\label{lem:crit}
With $q$ as above, $a$ satisfies \eqref{eq:rec} for all large $n$ if and only if
$u^{\!\top}M^{\,j}q(M)w=0$ for every $j\ge0$, and it is enough to check $j<S$.
\end{lemma}

\begin{proof}
Substituting the walk expression, the residual $a(n)-\sum_ic_ia(n-i)$ equals
$u^{\!\top}M^{\,g(n)-r}q(M)w$ up to the fixed linear reparametrisation $g$. For the bound,
$M^{S}$ is an integer combination of $I,M,\dots,M^{S-1}$ by Cayley--Hamilton, so the values
$u^{\!\top}M^{j}q(M)w$ satisfy the monic recurrence given by the characteristic polynomial of
$M$; $S$ consecutive zeros therefore force all later ones, and the last nonzero value pins the
threshold exactly.
\end{proof}

\begin{theorem}
$a(n)=f(n)$ for every $n\ge19$.
\end{theorem}

\begin{proof}
The vector $w'=q(M)w$ was formed by $48$ matrix--vector products in exact integer
arithmetic, and $u^{\!\top}M^{j}w'$ evaluated for $j=0,1,\dots$, again exactly. Those integers
vanish from the point corresponding to $n=67$ onwards, and the run was continued until the
working vector was identically zero, which settles every larger $j$ at once. So $a$ satisfies
\eqref{eq:rec} for every $n>66$. By Lemma~\ref{lem:ann} $f$ satisfies \eqref{eq:rec}
identically. Both were then evaluated in exact arithmetic at the $48$ consecutive indices
$n=67,\dots,114$ and agree at each. A monic recurrence of order $48$
determines $a(m)$ from $a(m-1),\dots,a(m-48)$, and for every $m\ge115$ both
$m>66$ and $m-48\ge67$ hold, so induction on $m$ gives $a(m)=f(m)$ for every
$m\ge67$.
Below $n=67$ the recurrence argument gives nothing, since \eqref{eq:rec} is only
established for $a$ from $n=67$ on. The remaining indices $n=19,\dots,66$ were
therefore checked one at a time, directly against the walk count in exact integer arithmetic,
and agree.
\end{proof}

The range is tight in the sense that was checked: at $n=18$ the closed form and the sequence differ, so no larger range is available.

\section{Verification}

Four checks, each independent of the algebra above.

First, the walk model reproduces every one of the $19$ terms the entry publishes,
exactly, in integer arithmetic. This is what ties the model to the entry: the model is built
by reading the entry's English, and a misreading gives a different digraph and different
counts.

Second, the closed form was evaluated in exact rational arithmetic --- not floating point ---
at every published term from $n=19$ on, and agrees with the entry's own DATA at each.

Third, the annihilation test of Lemma~\ref{lem:crit} was carried out exactly, and the model
and the closed form were then compared directly at sixty consecutive indices beyond
$n=19$, far past where the proof already settles the matter.

Fourth, the closed form was deliberately perturbed --- by a constant and by a term linear in
$n$ --- and each perturbed form was rejected by the same comparison, so the test is not one
that anything would pass.

\begin{thebibliography}{9}
\bibitem{oeis} The OEIS Foundation, \emph{The On-Line Encyclopedia of Integer
Sequences}, \url{https://oeis.org}, 2026. Sequence A227255.
\bibitem{stanley} R.~P.~Stanley, \emph{Enumerative Combinatorics}, Volume 1, 2nd ed.,
Cambridge University Press, 2011. (Transfer-matrix method, Section 4.7; $C$-finite sequences,
Section 4.1.)
\bibitem{kp} M.~Kauers and P.~Paule, \emph{The Concrete Tetrahedron}, Springer, 2011.
(Closed forms and linear recurrences with constant coefficients.)
\bibitem{horn} R.~A.~Horn and C.~R.~Johnson, \emph{Matrix Analysis}, 2nd ed., Cambridge
University Press, 2013.
\end{thebibliography}

\end{document}
